\section{Fase 4 -- \emph{Stack} de seguridad y CMS vulnerables}

\subsection{Objetivo}

Sobre el clúster con monitorización ya activa se despliegan dos detectores y tres aplicaciones vulnerables. Los detectores son \textbf{Falco} (HIDS basado en \texttt{eBPF}, observa \emph{syscalls} dentro de cualquier \emph{container}) y \textbf{Suricata} (NIDS basado en \texttt{af\_packet}, observa tráfico de red en \texttt{eth0} de un nodo). Las aplicaciones vulnerables son \textbf{DVWA}, \textbf{WebGoat} y \textbf{Juice Shop} --- tres CMS clásicos del catálogo OWASP que servirán como blancos de \emph{pen-testing} simulado desde una Kali externa al clúster.

Toda la fase se automatiza desde \texttt{scripts/setup-security.sh}, que comparte el patrón de \texttt{setup-lab.sh}: \emph{re-exec} como \texttt{master}, \texttt{helm upgrade -{}-install -{}-wait}, idempotencia por diseño.

\subsection{Arquitectura}

\begin{center}
\begin{tabular}{|l|l|l|l|}
\hline
\textbf{Componente} & \textbf{Despliegue} & \textbf{Nodo} & \textbf{Acceso} \\
\hline
Falco       & \texttt{DaemonSet} (sin master)         & worker1, worker2 & --- \\
Suricata    & \texttt{DaemonSet} (\texttt{nodeSelector}) & worker1       & \texttt{hostNetwork} eth0 \\
DVWA        & \texttt{Deployment} 1 \emph{replica}    & worker1          & NodePort 30430 \\
WebGoat     & \texttt{Deployment} 1 \emph{replica}    & worker1          & NodePort 30380 \\
Juice Shop  & \texttt{Deployment} 1 \emph{replica}    & worker1          & NodePort 31592 \\
\hline
\end{tabular}
\end{center}

Pipeline de eventos:
\begin{itemize}
  \item \textbf{Falco}: \texttt{stdout} JSON $\to$ Promtail (ya desplegado) $\to$ Loki $\to$ Grafana.
  \item \textbf{Suricata}: \texttt{eve.json} (archivo) $\to$ \emph{sidecar} \texttt{tail -F} $\to$ \texttt{stdout} $\to$ Promtail $\to$ Loki $\to$ Grafana.
  \item \textbf{Tráfico atacante}: Kali (10.0.1.20, externa al cluster CIDR) $\to$ \texttt{eth0} worker1 $\to$ visto por Suricata (\texttt{hostNetwork}) y por los \emph{logs} de Apache/Spring/Node de cada CMS.
\end{itemize}

\subsection{Concentración en worker1}

Los tres CMS y Suricata corren en el mismo nodo. Razón técnica: Suricata captura tráfico con \texttt{hostNetwork: true} desde \texttt{eth0} del nodo donde corre. Si DVWA migrara a otro nodo (por re-scheduling tras una caída), Suricata dejaría de ver el tráfico atacante. Concentrar todo en worker1 es la garantía simple de que las firmas siempre disparen.

Trade-off: worker1 (3072\,MiB de RAM) opera al 80--90\% de uso con todo el \emph{stack} activo. Los \texttt{resources.limits} de cada componente están ajustados para minimizar \emph{OOMKill} bajo carga:

\begin{itemize}
  \item Falco: 512\,Mi limit (default 1\,Gi).
  \item Suricata: 1\,Gi limit (default 256\,Mi --- insuficiente con ET Open completo cargado).
  \item DVWA: 256\,Mi limit (MySQL embebido es el componente más pesado).
  \item WebGoat: 512\,Mi limit (Spring Boot + heap acotado por \texttt{-Xmx384m}).
  \item Juice Shop: 192\,Mi limit (Node + Express + sqlite efímera, austero).
\end{itemize}

\subsection{Falco}

Falco se instala con el \emph{chart} oficial \texttt{falcosecurity/falco} (8.0.2, app 0.43.1).

\subsubsection{Driver \texttt{modern\_ebpf}}

\begin{lstlisting}[style=yaml]
driver:
  kind: modern_ebpf
\end{lstlisting}

\texttt{modern\_ebpf} usa \texttt{BPF\_PROG\_TYPE\_RAW\_TRACEPOINT}, lo que requiere kernel \texttt{\textgreater= 5.8}. Ubuntu 22.04 cumple (kernel 5.15). Ventajas frente al \emph{kernel module} tradicional o al \emph{eBPF legacy probe}:

\begin{itemize}
  \item Sin \texttt{kernel-headers} en los \emph{workers}: no hay compilación local de un módulo.
  \item Sin dependencia de \texttt{dkms}: el \emph{upgrade} del kernel no rompe la instalación de Falco.
  \item \emph{Self-contained}: el binario eBPF se compila \emph{ahead-of-time} y se carga vía syscall \texttt{bpf()}.
\end{itemize}

\subsubsection{Salida JSON a \texttt{stdout}}

\begin{lstlisting}[style=yaml]
falco:
  json_output: true
  json_include_output_property: true
  json_include_tags_property: true
  log_stderr: true
  log_syslog: false
  stdout_output:    { enabled: true }
  file_output:      { enabled: false }
  http_output:      { enabled: false }
\end{lstlisting}

Solo \texttt{stdout}, que kubelet escribe en \texttt{/var/log/pods/.../falco/0.log} $\to$ Promtail $\to$ Loki. \texttt{file\_output} saturaría \texttt{hostPath} sin aportar nada (los eventos ya están en Loki); \texttt{http\_output} requeriría un \emph{receiver} externo (\emph{falcosidekick}, que dejamos desactivado por la misma razón).

\subsubsection{El parámetro \texttt{base\_syscalls.repair}}

\begin{lstlisting}[style=yaml]
falco:
  base_syscalls:
    repair: true
\end{lstlisting}

Sin esta línea, los \emph{rules built-in} que filtran por \texttt{proc.tty != 0} (como \texttt{Terminal shell in container}) \textbf{nunca disparan}, porque \texttt{proc.tty} se queda siempre en 0 --- el árbol de procesos no tiene los datos para poblarlo. Es un \emph{footgun} silencioso: el \emph{ruleset} carga, los eventos se ven en \texttt{kubectl logs}, pero las reglas \emph{built-in} críticas no disparan nunca.

Con \texttt{repair: true} Falco analiza el \emph{ruleset}, infiere qué \emph{syscalls} auxiliares necesita (\texttt{clone}, \texttt{fork}, \texttt{setsid}, \texttt{procexit}, \texttt{ioctl}\dots) para poblar \texttt{proc.tty}, \texttt{proc.pname}, \texttt{proc.cmdline}, \texttt{container.id}, etc., y se suscribe automáticamente.

\subsubsection{Regla custom: ``Lab --- Shell en container''}

\begin{lstlisting}[style=yaml]
customRules:
  rules-lab.yaml: |-
    - rule: Lab -- Shell en container
      desc: >
        Dispara cuando un shell (bash, sh, zsh, ...) se ejecuta dentro
        de un container. Excluye casos legitimos donde un package manager
        o el runtime spawnea shells como hooks de instalacion.
      condition: >
        spawned_process and container
        and proc.name in (bash, sh, zsh, ksh, csh, dash, ash, fish)
        and not proc.pname in (apt, apt-get, dpkg, dpkg-reconfigure, yum,
                               dnf, apk, pip, npm, conda, sudo)
      output: >
        [LAB] Shell en container
        (pod=%k8s.pod.name ns=%k8s.ns.name container=%container.name
         shell=%proc.name cmdline=%proc.cmdline parent=%proc.pname tty=%proc.tty)
      priority: WARNING
      tags: [lab, shell]
\end{lstlisting}

\paragraph{Iteración hasta la regla final.} La versión inicial (\emph{canary} de validación E2E del \emph{probe}) era \texttt{spawned\_process and container}, sin filtros adicionales. Funcionaba como prueba de vida: cualquier \texttt{execve} dentro de cualquier \emph{container} disparaba un evento. El problema observado:

\begin{itemize}
  \item Flannel ejecuta \texttt{iptables} y \texttt{ip} continuamente.
  \item \texttt{kube-proxy} sondea reglas y servicios cada pocos segundos.
  \item Los \emph{hooks} de \texttt{apt}/\texttt{dpkg} durante actualizaciones automáticas lanzan \emph{shells} legítimos.
\end{itemize}

Resultado: $\sim$16 eventos por segundo en estado \emph{idle}, casi todos ruido del \emph{runtime}. En 10 minutos, $\sim$9\,656 eventos --- imposible distinguir un atacante real.

\paragraph{Refinamiento.} Se filtró por \texttt{proc.name in (bash, sh, zsh, \dots)} para quedarse con \emph{shells} interactivas, y se excluyó \texttt{proc.pname in (apt, dpkg, sudo, \dots)} para descartar \emph{shells} legítimas como hijos de gestores de paquetes. La regla refinada genera 0 eventos en estado \emph{idle} y dispara puntualmente cuando un humano hace \texttt{kubectl exec -- sh}.

\subsubsection{\texttt{customRules} y \texttt{rules\_files}: la trampa}

\begin{lstlisting}[style=yaml]
falco:
  rules_files:
    - /etc/falco/falco_rules.yaml
    - /etc/falco/falco_rules.local.yaml
    - /etc/falco/rules.d
\end{lstlisting}

El \emph{chart} monta \texttt{customRules} como \texttt{ConfigMap} en \texttt{/etc/falco/rules.d/}. \textbf{Crítico}: \texttt{/etc/falco/rules.d} debe estar en \texttt{rules\_files}; si se sobrescribe \texttt{rules\_files} sin incluirla, las reglas \emph{custom} se montan en el pod pero Falco \emph{nunca las carga}, sin error visible. Es la misma categoría de fallo que el \texttt{base\_syscalls.repair}: el síntoma es ``mi regla no dispara'' y la causa está enterrada en una sub-clave de configuración.

\subsection{Suricata}

Suricata no tiene \emph{chart} oficial estable. Se desarrolló un \emph{chart} propio en \texttt{charts/suricata/} con tres recursos: \texttt{ConfigMap} (\texttt{suricata.yaml} + \texttt{update.yaml}), \texttt{DaemonSet} (init + main + sidecar) y los \emph{helpers} de Helm.

\subsubsection{\texttt{hostNetwork} y \emph{capabilities}}

\begin{lstlisting}[style=yaml]
hostNetwork: true
dnsPolicy: ClusterFirstWithHostNet
securityContext:
  capabilities:
    add: [NET_ADMIN, NET_RAW, SYS_NICE]
\end{lstlisting}

\texttt{hostNetwork: true} es el \emph{toggle} crítico: el pod comparte el \emph{network namespace} del nodo y ve \texttt{eth0} entera. Sin esto, Suricata solo ve el \emph{netns} vacío del pod y captura 0 paquetes útiles.

Las \emph{capabilities} mínimas son:
\begin{itemize}
  \item \texttt{NET\_ADMIN}: setear modo promiscuo, configurar el \emph{ring buffer} \texttt{TPACKET\_V3}.
  \item \texttt{NET\_RAW}: crear \emph{sockets} \texttt{AF\_PACKET} (capa 2, Ethernet completo).
  \item \texttt{SYS\_NICE}: cambiar prioridad de \emph{threads} de captura.
\end{itemize}

Notablemente, \textbf{no se usa} \texttt{privileged: true}. El kernel Linux permite todo lo necesario con \emph{capabilities} granulares, sin dar acceso total al host.

\paragraph{\texttt{dnsPolicy: ClusterFirstWithHostNet}.} Cuando \texttt{hostNetwork: true}, el pod hereda \texttt{/etc/resolv.conf} del nodo (que no apunta a CoreDNS). Si tuviéramos que resolver \emph{services} de Kubernetes desde dentro del pod, fallarían. Con \texttt{ClusterFirstWithHostNet} se fuerza el uso de CoreDNS aunque la red sea del host. En el lab Suricata no resuelve nombres K8s, pero la línea queda como higiene.

\subsubsection{Captura con \texttt{af-packet}}

\begin{lstlisting}[style=yaml]
af-packet:
  - interface: eth0
    cluster-id: 99
    cluster-type: cluster_flow
    defrag: yes
    use-mmap: yes        # ring buffer en memoria compartida
    tpacket-v3: yes      # API moderna del kernel
\end{lstlisting}

Modo IDS pasivo (no IPS). Para IPS habría que usar \texttt{nfqueue} y enganchar el tráfico vía \texttt{iptables -j NFQUEUE}, lo que no funciona bien en \emph{containers} y queda fuera de \emph{scope} del lab.

\subsubsection{Variables de red: \texttt{HOME\_NET} y \texttt{HTTP\_PORTS}}

\begin{lstlisting}[style=yaml]
vars:
  address-groups:
    HOME_NET: "[10.0.1.10/32,10.0.1.11/32,10.0.1.12/32,10.233.0.0/16]"
    EXTERNAL_NET: "!$HOME_NET"
  port-groups:
    HTTP_PORTS: "80,30430,30380,31592"
\end{lstlisting}

\paragraph{Diseño del \texttt{HOME\_NET}.} Las firmas ET Open están escritas en términos de \texttt{\$EXTERNAL\_NET -> \$HOME\_NET}: solo disparan cuando origen y destino tienen ese sentido. Si la IP de la Kali atacante (10.0.1.20) cayese dentro de \texttt{HOME\_NET}, las firmas \emph{no dispararían} (sería \texttt{HOME -> HOME}). Por eso se enumeran nodos individuales (\texttt{/32}) y el CIDR de pods/services (\texttt{10.233.0.0/16}) explícitamente, dejando 10.0.1.20 fuera. Era el \emph{footgun} clásico: usar \texttt{10.0.1.0/24} ``por simplicidad'' anula la mayoría de detecciones sin error visible.

\paragraph{\texttt{HTTP\_PORTS} parametrizado.} Las firmas \texttt{ET WEB\_SERVER \dots} están escritas como \texttt{flow:to\_server,established; http.uri; \dots}, lo que requiere que Suricata haya clasificado el flujo como HTTP \emph{antes} de evaluar la regla. La clasificación se hace por \texttt{\$HTTP\_PORTS}. Si un CMS expone en un puerto que no está en esa lista, Suricata lo trata como TCP genérico y las firmas web nunca disparan --- de nuevo, sin error visible.

Cuando se añadieron WebGoat (30380) y Juice Shop (31592), tuvieron que añadirse a \texttt{HTTP\_PORTS} para que las firmas web disparasen contra ellos. Para no repetir el error, la lista vive como variable en \texttt{values.yaml} del \emph{chart}:

\begin{lstlisting}[style=yaml]
suricata:
  httpPorts:
    - 80
    - 30430
    - 30380
    - 31592
\end{lstlisting}

y el \texttt{ConfigMap} la materializa con la función \texttt{join} de Sprig:

\begin{lstlisting}[style=yaml]
HTTP_PORTS: "{{ join "," .Values.suricata.httpPorts }}"
\end{lstlisting}

\subsubsection{El \emph{init container}: descarga del \emph{ruleset} ET Open}

\begin{lstlisting}[style=yaml]
initContainers:
  - name: rules-update
    image: jasonish/suricata:latest
    command: ["suricata-update"]
    args:
      - "--no-test"
      - "--output=/var/lib/suricata/rules"
      - "--data-dir=/var/lib/suricata"
      - "--config=/etc/suricata/update.yaml"
\end{lstlisting}

\texttt{suricata-update} viene en la imagen oficial. El \emph{ruleset} se descarga a un \texttt{emptyDir} compartido con el \emph{container} principal (no se persiste --- se re-descarga en cada \emph{start} del pod, $\sim$30\,s).

\paragraph{URL directa frente al alias \texttt{et/open}.}

\begin{lstlisting}[style=yaml]
update.yaml: |
  sources:
    - https://rules.emergingthreats.net/open/suricata-8.0.0/emerging.rules.tar.gz
\end{lstlisting}

El alias \texttt{et/open} requiere que \texttt{suricata-update update-sources} haya cargado el catálogo de fuentes \emph{antes} del \emph{fetch}. Ese catálogo se guarda en \texttt{-{}-data-dir}, que en nuestro setup es un \texttt{emptyDir} nuevo cada arranque. Como el \emph{init container} no ejecuta \texttt{update-sources}, el alias no resuelve --- se intenta interpretar como URL literal y falla con \texttt{unknown url type: 'et/open'}. Pero el \emph{init container} sale con \texttt{exit 0} (\emph{warning}, no error), así que el pod queda \texttt{Ready} con solo 436 reglas \emph{decoder/protocol} y \textbf{cero firmas ET Open}. Detección silenciosamente rota. Fijar la URL directa elimina el problema.

\subsubsection{El \emph{sidecar} \texttt{tail -F eve.json}}

\begin{lstlisting}[style=yaml]
- name: log-tail
  image: jasonish/suricata:latest
  command: ["sh", "-c"]
  args:
    - |
      until [ -f /var/log/suricata/eve.json ]; do sleep 1; done
      exec tail -F /var/log/suricata/eve.json
  securityContext:
    allowPrivilegeEscalation: false
    capabilities: { drop: [ALL] }
\end{lstlisting}

Promtail solo lee \texttt{stdout} de los pods; Suricata escribe \texttt{eve.json} a archivo. El \emph{sidecar} es el puente: \texttt{tail -F} (mayúscula) re-abre el archivo si Suricata lo rota (algo que \texttt{tail -f} minúscula \emph{no} hace, se queda leyendo el \emph{inode} antiguo eternamente).

\paragraph{Reuso de la imagen principal.} El \emph{sidecar} usa \texttt{jasonish/suricata} en vez de \texttt{busybox}. Razones: (1) \texttt{tail} GNU está en la imagen, (2) un solo \emph{pull} y un solo conjunto de \emph{layers} cacheados, (3) algunos \emph{blobs} de \texttt{busybox} en Docker Hub se sirven desde Cloudflare R2 con IPs ocasionalmente no enrutables en la red \texttt{red-prueba}. Trade-off: el \emph{sidecar} arranca con 139\,MB de imagen en vez de 5\,MB, pero no añade RAM real (\emph{layers} compartidos).

\paragraph{Defensa en profundidad.} El \emph{sidecar} solo lee --- no necesita \emph{capabilities} ni acceso de escritura. \texttt{drop: [ALL]} elimina cualquier privilegio remanente. Un compromiso del \emph{sidecar} no puede manipular \emph{logs} ni capturar tráfico.

\subsubsection{Coherencia config: \texttt{stats}}

\begin{lstlisting}[style=yaml]
outputs:
  - eve-log:
      enabled: yes
      types:
        - alert
        - http
        - dns
        - tls
        - flow
        # NO incluir 'stats' aqui si outputs.stats.enabled = no
  - stats:
      enabled: no
\end{lstlisting}

Una particularidad de Suricata: si \texttt{outputs.stats.enabled: no} globalmente, no se puede incluir \texttt{stats} dentro de \texttt{eve-log.types}. Suricata aborta el arranque con \emph{``stats are disabled globally''}. La coherencia entre el \emph{sub-output} y el \emph{output} principal es obligatoria.

\subsection{DVWA --- \emph{Damn Vulnerable Web Application}}

DVWA se despliega vía \emph{chart} propio en \texttt{charts/dvwa/}.

\subsubsection{Imagen \texttt{vulnerables/web-dvwa}}

\begin{lstlisting}[style=yaml]
image:
  repository: vulnerables/web-dvwa
  tag: latest
\end{lstlisting}

Se eligió la imagen \texttt{vulnerables/web-dvwa} (todo-en-uno: Apache + PHP + MySQL embebido) en lugar de \texttt{ghcr.io/digininja/dvwa} (la oficial, que requiere MariaDB como segundo \texttt{Deployment}). Razón: un solo \emph{container}, menos partes móviles, BD inicializada \emph{out of the box}.

\subsubsection{Sin persistencia (decisión consciente)}

\begin{lstlisting}[style=yaml]
persistence:
  enabled: false
\end{lstlisting}

La imagen \texttt{vulnerables/web-dvwa} trae la base de datos MariaDB \textbf{horneada en la imagen} (en \texttt{/var/lib/mysql}). Su \emph{entrypoint} (\texttt{/main.sh}) solo hace \texttt{service mysql start}, no ejecuta \texttt{mysql\_install\_db}. Si montamos un PVC encima de \texttt{/var/lib/mysql}, el \emph{mount} ``ensombrece'' la BD pre-cocinada $\to$ \emph{datadir} vacío $\to$ \texttt{mysqld} arranca pero falla con \emph{``Table 'mysql.user' doesn't exist''}.

Alternativas descartadas: pre-popular el PV con un \emph{init container} que copie la BD baked, o ejecutar \texttt{mysql\_install\_db} manualmente. Ambas añaden complejidad sin beneficio práctico --- en un lab de detección de ataques, perder progreso entre \emph{restarts} es aceptable.

\subsubsection{Probe \texttt{/login.php}}

\begin{lstlisting}[style=yaml]
probes:
  liveness:
    path: /login.php
    initialDelaySeconds: 30
    periodSeconds: 30
\end{lstlisting}

El \emph{probe} apunta a \texttt{/login.php}, no a \texttt{/}. Razón: \texttt{/} devuelve 302 redirigiendo a \texttt{/setup.php} hasta que la BD se inicializa. \texttt{/login.php} devuelve 200 directamente con la BD horneada operativa al arranque.

\subsubsection{\texttt{strategy: Recreate}}

\begin{lstlisting}[style=yaml]
strategy:
  type: Recreate
\end{lstlisting}

Aunque no hay PVC, se mantiene el patrón por consistencia con cualquier futuro \emph{enable} de \texttt{persistence}. Con \texttt{Recreate}, el pod viejo se mata antes de arrancar el nuevo --- no hay dos \texttt{mysqld} escribiendo a la misma BD simultáneamente.

\subsection{WebGoat}

\subsubsection{Imagen oficial}

\begin{lstlisting}[style=yaml]
image:
  repository: webgoat/webgoat
  tag: latest
\end{lstlisting}

\texttt{webgoat/webgoat} es la imagen oficial OWASP. Lleva WebGoat (puerto 8080) y WebWolf (puerto 9090) en un único \emph{container}, lanzados por el \emph{entrypoint} del JAR.

\subsubsection{Variables de entorno críticas}

\begin{lstlisting}[style=yaml]
env:
  - { name: WEBGOAT_HOST, value: "0.0.0.0" }
  - { name: WEBGOAT_PORT, value: "8080" }
  - { name: WEBWOLF_HOST, value: "0.0.0.0" }
  - { name: WEBWOLF_PORT, value: "9090" }
  - { name: JAVA_TOOL_OPTIONS, value: "-Xmx384m -Xms256m" }
\end{lstlisting}

\paragraph{\texttt{WEBGOAT\_HOST=0.0.0.0}.} Sin esto, el JAR bindea solo a \texttt{127.0.0.1} y el \texttt{NodePort} responde con \emph{``Connection refused''}: el tráfico llega al \emph{container} pero no a ningún socket abierto. Es un \emph{default} desafortunado de la imagen oficial.

\paragraph{\texttt{JAVA\_TOOL\_OPTIONS} para acotar el \emph{heap}.} Sin acotar, Spring Boot crece hasta donde le deja el host --- en worker1, eso significa rebasar el limit de 512\,Mi del \emph{container} y morir \texttt{OOMKilled} (\emph{exit 137}) tras pocos minutos. Acotar \texttt{-Xmx384m} reserva $\sim$128\,Mi para \emph{metaspace}, \emph{thread stacks} y código JIT. Cualquier JVM moderna lee \texttt{JAVA\_TOOL\_OPTIONS} del entorno automáticamente.

\subsubsection{URL \emph{case-sensitive}: \texttt{/WebGoat}}

WebGoat usa \texttt{/WebGoat} con \texttt{W} y \texttt{G} mayúsculas en sus redirecciones internas. \texttt{/webgoat} devuelve 404. El \emph{probe} canónico:

\begin{lstlisting}[style=yaml]
probes:
  liveness:
    path: /WebGoat/login.mvc
    initialDelaySeconds: 120
\end{lstlisting}

\texttt{initialDelaySeconds: 120} cubre el \emph{cold start} de Spring Boot (\emph{classpath scanning} + inicialización de H2 embedded), que puede llegar a 90\,s en worker1 cargado.

\subsection{Juice Shop}

\subsubsection{Despliegue mínimo}

\begin{lstlisting}[style=yaml]
image:
  repository: bkimminich/juice-shop
  tag: latest
service:
  type: NodePort
  port: 3000
  nodePort: 31592
persistence:
  enabled: false
\end{lstlisting}

Juice Shop es el más austero del trío: Node + Express + sqlite efímera. No requiere variables de entorno específicas; la imagen oficial corre en modo \emph{production} por defecto y los retos están todos accesibles desde \texttt{/\#/score-board}.

\subsubsection{Sin login inicial}

A diferencia de DVWA (\texttt{admin/password}) y WebGoat (registro de cuenta nueva en cada arranque), Juice Shop expone toda la \emph{landing page} sin autenticación. Los retos OWASP Top 10 están integrados como minijuegos progresivos.

\subsection{El \emph{script} \texttt{setup-security.sh}}

Estructura en 9 fases:

\begin{enumerate}
  \item \textbf{Precondición}: verifica que \texttt{monitoring/promtail} esté \texttt{Ready}. Sin Promtail, los eventos generados por Falco y Suricata se perderían --- mejor abortar pronto que dejar detectores ``huérfanos''.
  \item \textbf{Repo Helm de Falco}: solo Falco viene de un repo \emph{upstream}; Suricata, DVWA, WebGoat y Juice Shop son \emph{charts} locales bajo \texttt{charts/}.
  \item \textbf{Namespaces}: \texttt{security} + \texttt{vulnerable} con el patrón idempotente \texttt{create -{}-dry-run=client -o yaml | apply -f -}.
  \item \textbf{Falco}.
  \item \textbf{Suricata}: \texttt{helm install /vagrant/charts/suricata}.
  \item \textbf{DVWA}: \texttt{helm install /vagrant/charts/dvwa}.
  \item \textbf{WebGoat}: \texttt{-{}-timeout 8m} para cubrir el \emph{cold start} de Spring Boot.
  \item \textbf{Juice Shop}.
  \item \textbf{Resumen}: imprime endpoints y queries LogQL canónicas.
\end{enumerate}

Invocación:

\begin{lstlisting}[style=bash]
vagrant provision test-master --provision-with setup-security
\end{lstlisting}

\subsection{Lecciones transversales}

\subsubsection{Charts custom: límite a Helm \emph{upstream}}

Cuando no hay \emph{chart} estable y bien mantenido (Suricata) o cuando el oficial impone más complejidad que beneficio (DVWA con la oficial \texttt{digininja/dvwa} + MariaDB aparte), escribir un \emph{chart} propio simplificado de 3--4 plantillas es \emph{más} robusto que parchear un \emph{chart} \emph{upstream} con \texttt{values}. El coste de mantenimiento es contenido: los \emph{templates} son lineales y los \texttt{values} viven con el código que documentan.

\subsubsection{Imágenes con \emph{defaults} hostiles a \texttt{NodePort}}

WebGoat es el ejemplo claro: una imagen ``oficial'' que bindea a \texttt{127.0.0.1} por defecto es prácticamente inservible en Kubernetes sin un \texttt{env} explícito de \texttt{0.0.0.0}. La lección: nunca asumir que una imagen oficial está optimizada para el caso \emph{containerizado} --- probar el \emph{flow} \texttt{NodePort $\to$ pod} pronto y observar logs.

\subsubsection{Detectores silenciosamente rotos}

Tres ejemplos en este capítulo:
\begin{itemize}
  \item Falco sin \texttt{base\_syscalls.repair} $\to$ \emph{rules built-in} con \texttt{proc.tty} nunca disparan.
  \item Suricata con alias \texttt{et/open} en \texttt{update.yaml} $\to$ \emph{ruleset} ET Open no se descarga, queda con 436 \emph{decoder rules}.
  \item Suricata sin \texttt{HTTP\_PORTS} actualizado $\to$ firmas \texttt{ET WEB\_SERVER} no clasifican el flujo como HTTP y no disparan.
\end{itemize}

Los tres son síntomas de la misma categoría: ``el sistema funciona aparentemente, los pods están \texttt{Ready}, los \emph{logs} no muestran errores, pero los detectores están sordos''. La defensa es el test \emph{end-to-end}: lanzar un ataque conocido (un \texttt{curl} con \texttt{?id=1' OR '1'='1}, un \texttt{kubectl exec -- sh}) y verificar que aparece el evento esperado en Loki/Grafana \emph{antes} de dar el \emph{stack} por terminado.

\subsection{Estado tras Fase 4}

Con las cuatro fases completas, un \texttt{vagrant up} seguido de tres \texttt{vagrant provision -{}-provision-with} deja un laboratorio Kubernetes funcional, observable y con superficie de ataque definida, listo para \emph{pen-test} simulado desde Kali. Los siguientes hitos no cubiertos en este documento (y reservados para iteraciones posteriores) incluyen el cableado de Alertmanager y un \emph{bot} de Telegram para alertas \emph{push}.
